\scanpage{3}

One way of seeing it is to say that
the datum of $\uline{s}$ defines a splitting
\inserted{$\varphi_{\uline{s}}$} of
the extension $(*)$ above, and that \unsure{viewed as}
a subgroup of $\operatorname{SRT}(X)$,
$S_{\unsure{1|*}}T(X,\uline{s})$ is identified
with the inverse image of
\[
\varphi_{\uline{s}}(T(X))\subset\operatorname{SUT}(X)
\]
under the canonical homomorphism
\[
\operatorname{SUR}(X)\longrightarrow\operatorname{SUT}(X)
\]
of passage to the quotient. The change from
$\uline{s}$ to $\uline{s}'=u.\uline{s}$
transforms $\varphi_{\uline{s}}$ into
\[
\varphi_{\uline{s}'}=\operatorname{int}(u)\circ\varphi_{\uline{s}}
\]
thus (if $\widetilde u\in\mathbb R^I$ lifts $u$) one finds
\[
S_{\unsure{1|*}}T(X,\uline{s}')
=\operatorname{int}(\widetilde u)S_{\unsure{1|*}}T(X,\uline{s}).
\]

In terms of $\vec X$, the definition of
$S_{\unsure{1|*}}T(X,\uline{s})$
depends on a choice $\uline{s}$ that is \uline{not innocuous},
on the other hand that of $S_1RT(X)$ depends only on
an innocuous choice, that of the family of
toric structures $(t_i)_{i\in I}$.
Moreover $\prod T_i$ as
a torsor under $\unsure{\mathbb U^I}$ is realized
\inserted{(not faithfully, incidentally)}
as \unsure{a class} of splittings of the extension $(*)$
(via inner automorphisms by $\mathbb U^I$)
$\bigl[$\unsure{deduced} from $\operatorname{SRT}(X)$
with \unsure{its} additional structures$\bigr\}$
namely\unsure{)|,} the inclusions

\diagram{
\begin{tikzcd}[row sep=small]
\mathbb R^I\arrow[dr] &\\
&\operatorname{SRT}(X)\\
S_{\unsure{1|*}}T^!(X)\arrow[ur]&
\end{tikzcd}
}

\unsure{coinciding|coinciding} on $\mathbb Z^I$.

\margin{2}
